% Prose rendering of PrimeTensor/Order.lean.
% Fragment: no \documentclass, no preamble, no document environment.

\section{Native order and multiplicative annuli}
\label{Order:sec}

\leanfile{PrimeTensor/Order.lean}

\subsection*{What this file establishes}

Two earlier modules left a gap that this one closes.
\texttt{PrimeTensor/Ratio.lean} defined the cross-product strict order
$\name{Lt}$ on representatives and proved nothing about it;
\texttt{PrimeTensor/MulRat.lean} took the quotient and transported only the
multiplicative structure across, leaving $\name{MulRat}$ with no order at all.
Here $\name{Lt}$ is shown to respect the equivalence, lifted to
$\name{MulRat}$, and proved irreflexive and transitive.

On top of that comes the notion the completion will actually be built from:
$\name{Within}(\delta, a, b)$, saying that $a$ and $b$ lie in a common
multiplicative $\delta$-annulus.  There is no distance function anywhere, and
no subtraction --- closeness is expressed by bounding a ratio in both
directions, and the tolerance is a magnitude exceeding $1$ rather than a small
positive number.

\subsection{Positivity, again}

\begin{lemma}[\lean{PrimeRatio.multiset\_eval\_pos\_order}]\label{Order:lem:pos}
$0 < \name{eval}(q)$ for every $q : \name{PrimeMultiset}$.
\end{lemma}

\begin{proof}
By $\name{Quotient.inductionOn}$, reducing to a representative bag, where
$0 < 1$ and $\name{PrimeBag.one\_le\_eval}$ compose.
\end{proof}

This is word for word the lemma proved in
\texttt{PrimeTensor/MulRat.lean}, restated because that one is \lean{private}
and so not visible outside its own file.  Both are \lean{private}, and the
names differ only by the suffix $\lean{\_order}$.  Everything below rests on
it: each order proof cancels a factor, and cancellation of a strict inequality
on $\mathbb{N}$ is valid only against a positive factor.

\subsection{The order respects the equivalence}

\begin{lemma}[Left congruence, \lean{PrimeRatio.lt\_same\_left}]
\label{Order:lem:left}
If $\name{Same}(a, a')$ then $\name{Lt}(a,b) \leftrightarrow \name{Lt}(a',b)$.
\end{lemma}

\begin{leancode}
private theorem lt_of_same_left {a a' b : PrimeRatio}
    (ha : Same a a') (h : Lt a b) : Lt a' b := by
  unfold Same at ha
  unfold Lt at h ⊢
  exact (Nat.mul_lt_mul_right (multiset_eval_pos_order a.lower)).mp (by
    calc
      (a'.upper.eval * b.lower.eval) * a.lower.eval
          = (a'.upper.eval * a.lower.eval) * b.lower.eval := by ac_rfl
      _ = (a.upper.eval * a'.lower.eval) * b.lower.eval := by rw [← ha]
      _ = (a.upper.eval * b.lower.eval) * a'.lower.eval := by ac_rfl
      _ < (b.upper.eval * a.lower.eval) * a'.lower.eval :=
          (Nat.mul_lt_mul_right (multiset_eval_pos_order a'.lower)).mpr h
      _ = (b.upper.eval * a'.lower.eval) * a.lower.eval := by ac_rfl)
\end{leancode}

\begin{proof}
Write $u_x = \name{eval}(x.\name{upper})$, $l_x = \name{eval}(x.\name{lower})$.
The forward direction is the private lemma \lean{lt\_of\_same\_left}
displayed above.  Its hypotheses are $u_a l_{a'} = u_{a'} l_a$ and
$u_a l_b < u_b l_a$, and its goal is $u_{a'} l_b < u_b l_{a'}$.

Since $\mathbb{N}$ has no division, the goal is multiplied by $l_a$ and the
factor cancelled at the end: by $\name{Nat.mul\_lt\_mul\_right}$ against the
positive $l_a$ it suffices to show
$(u_{a'} l_b)\,l_a < (u_b l_{a'})\,l_a$, which is the displayed chain ---
reassociate to expose $u_{a'} l_a$, rewrite it backwards along the hypothesis
to $u_a l_{a'}$, reassociate to expose $u_a l_b$, apply the strict hypothesis
multiplied by the positive $l_{a'}$, and reassociate once more.  Every
equality step is $\ctor{ac\_rfl}$.

The stated \lean{iff} then follows by applying the private lemma in both
directions, the second time to $\name{same\_symm}\,ha$.
\end{proof}

\begin{lemma}[Right congruence, \lean{PrimeRatio.lt\_same\_right}]
\label{Order:lem:right}
If $\name{Same}(b, b')$ then $\name{Lt}(a,b) \leftrightarrow \name{Lt}(a,b')$.
\end{lemma}

\begin{proof}
The mirror image of Lemma~\ref{Order:lem:left}, via the private
\lean{lt\_of\_same\_right}: multiply by $l_b$, and this time apply the strict
hypothesis \emph{first} and rewrite along $\name{Same}(b,b')$ afterwards,
since the hypothesis being replaced now sits on the right of the inequality.
\end{proof}

\begin{lemma}[\lean{PrimeRatio.lt\_iff\_of\_same}]\label{Order:lem:both}
If $\name{Same}(a,a')$ and $\name{Same}(b,b')$ then
$\name{Lt}(a,b) \leftrightarrow \name{Lt}(a',b')$.
\end{lemma}

\begin{proof}
Compose Lemmas~\ref{Order:lem:left} and~\ref{Order:lem:right}: replacing one
argument at a time is enough, and the two replacements are independent.
\end{proof}

\subsection{Irreflexivity and transitivity at representative level}

\begin{lemma}[\lean{lt\_irrefl\_ratio}]\label{Order:lem:irrefl-rep}
$\neg\, \name{Lt}(a,a)$ for every $a : \name{PrimeRatio}$.
\end{lemma}

\begin{proof}
Unfolded, $\name{Lt}(a,a)$ is $u_a l_a < u_a l_a$, contradicted by
irreflexivity of $<$ on $\mathbb{N}$.  Note that no positivity is needed
here --- the two sides are literally the same natural number.
\end{proof}

\begin{lemma}[\lean{lt\_trans\_ratio}]\label{Order:lem:trans-rep}
If $\name{Lt}(a,b)$ and $\name{Lt}(b,c)$ then $\name{Lt}(a,c)$.
\end{lemma}

\begin{proof}
The hypotheses are $u_a l_b < u_b l_a$ and $u_b l_c < u_c l_b$ and the goal is
$u_a l_c < u_c l_a$.  Multiplying the goal by the positive $l_b$ and cancelling
afterwards, it suffices to prove
$(u_a l_c)\,l_b < (u_c l_a)\,l_b$, and the chain runs
\[
  (u_a l_c) l_b = (u_a l_b) l_c
    < (u_b l_a) l_c = (u_b l_c) l_a
    < (u_c l_b) l_a = (u_c l_a) l_b,
\]
with the two strict steps being the hypotheses multiplied by the positive
$l_c$ and $l_a$ respectively, and the three equalities $\ctor{ac\_rfl}$.
\end{proof}

\begin{remark}[One pattern, three times]\label{Order:rem:pattern}
Lemmas~\ref{Order:lem:left}, \ref{Order:lem:right}
and~\ref{Order:lem:trans-rep} are the same manoeuvre: multiply through by the
denominator that the conclusion does not mention, walk a chain of
$\ctor{ac\_rfl}$ reassociations punctuated by the hypotheses, then cancel.  It
is the strict-inequality counterpart of the calculation proving transitivity of
$\name{Same}$ in \texttt{PrimeTensor/MulRat.lean}, and it exists for the same
reason: the natural formulation would divide, and $\mathbb{N}$ has no
division.  Each use needs a positive factor, and each gets it from
Lemma~\ref{Order:lem:pos}.
\end{remark}

\subsection{Lifting the order to the quotient}

\begin{definition}[\lean{MulRat.lt} and the \lean{LT} instance]
\label{Order:def:lt}
$\name{lt} : \name{MulRat} \to \name{MulRat} \to \mathrm{Prop}$ is
$\name{PrimeRatio.Lt}$ lifted through the quotient in both arguments, and
\lean{instance : LT MulRat} makes $a < b$ mean $\name{lt}(a,b)$.
\end{definition}

\begin{leancode}
def lt (a b : MulRat) : Prop :=
  Quotient.liftOn₂ a b
    PrimeRatio.Lt
    (by
      intro a₁ b₁ a₂ b₂ ha hb
      change PrimeRatio.Same a₁ a₂ at ha
      change PrimeRatio.Same b₁ b₂ at hb
      apply propext
      exact PrimeRatio.lt_iff_of_same ha hb)
\end{leancode}

\begin{remark}[This lift needs \lean{propext}]\label{Order:rem:propext}
The lifts in \texttt{PrimeTensor/MulRat.lean} produced elements of
$\name{MulRat}$, and their obligations were discharged by
$\name{Quotient.sound}$, which turns a $\name{Same}$ into an equality of
classes.  This one produces an element of $\mathrm{Prop}$, so its obligation is
that the two propositions are \emph{equal}, not merely equivalent.
Lemma~\ref{Order:lem:both} supplies an $\leftrightarrow$, and
$\name{propext}$ --- propositional extensionality --- converts it to the
required equality.  That axiom is the only extra principle the order costs,
and it is one Lean's standard library assumes throughout.
\end{remark}

\begin{theorem}[The lifted order is a strict order]\label{Order:thm:strict}
For $a, b, c : \name{MulRat}$: $\neg\, a < a$, and $a < b$ together with
$b < c$ gives $a < c$.  These are \lean{MulRat.lt\_irrefl} and
\lean{MulRat.lt\_trans}.
\end{theorem}

\begin{proof}
Both by $\name{Quotient.inductionOn}$, reducing to representatives, where the
goals become $\neg\,\name{Lt}(x,x)$ and $\name{Lt}(x,z)$ and the
representative-level Lemmas~\ref{Order:lem:irrefl-rep}
and~\ref{Order:lem:trans-rep} apply.  Transitivity inducts on all three
arguments in turn, carrying the hypotheses along.  No appeal to
$\name{Quotient.sound}$ is needed: the statements are propositions about
classes, not equations between them.
\end{proof}

\begin{remark}[What is not established]\label{Order:rem:order-gaps}
No order-theoretic instance is registered beyond \lean{LT}: there is no
\lean{Preorder}, \lean{PartialOrder} or \lean{LinearOrder} for
$\name{MulRat}$, so Mathlib's order lemmas do not apply and every fact must be
cited from this file.  Asymmetry is not stated, though it follows from
irreflexivity and transitivity.  Neither is trichotomy --- that of $a < b$,
$a = b$, $b < a$ exactly one holds --- nor any compatibility between the order
and the group structure, such as $a < b$ implying $ac < bc$.  Nothing relates
$<$ to $\name{inv}$.
\end{remark}

\subsection{Ratios and annuli}

\begin{definition}[\lean{MulRat.ratio}]\label{Order:def:ratio}
$\name{ratio}(a,b) \defeq a \cdot b^{-1}$.
\end{definition}

\begin{lemma}[\lean{MulRat.ratio\_self}]\label{Order:lem:ratio-self}
$\name{ratio}(a,a) = 1$.
\end{lemma}

\begin{proof}
Unfolding, this is $a \cdot a^{-1} = 1$, which is \lean{MulRat.mul\_inv} from
\texttt{PrimeTensor/MulRat.lean}.  It is tagged \lean{@[simp]}.
\end{proof}

This is a real strengthening over the representative level.  In
\texttt{PrimeTensor/Ratio.lean}, $\name{PrimeRatio.ratio}(a,a)$ is the term
$\langle u_a l_a,\, l_a u_a \rangle$, which is only $\name{Same}$ to the unit,
not equal to it.  After the quotient the two coincide, so the self-ratio can
be rewritten away by $\mathsf{simp}$ rather than carried as a relation.

\begin{definition}[Annulus membership, \lean{MulRat.Within}]
\label{Order:def:within}
For $\delta, a, b : \name{MulRat}$,
\[
  \name{Within}(\delta, a, b) \defeq
    \bigl(\name{ratio}(a,b) < \delta\bigr) \wedge
    \bigl(\name{ratio}(b,a) < \delta\bigr).
\]
\end{definition}

\begin{leancode}
def Within (δ a b : MulRat) : Prop :=
  ratio a b < δ ∧ ratio b a < δ
\end{leancode}

\begin{lemma}[\lean{within\_refl} and \lean{within\_symm}]
\label{Order:lem:within}
If $1 < \delta$ then $\name{Within}(\delta, a, a)$; and
$\name{Within}(\delta, a, b)$ implies $\name{Within}(\delta, b, a)$.
\end{lemma}

\begin{proof}
For reflexivity, Lemma~\ref{Order:lem:ratio-self} rewrites both conjuncts to
$1 < \delta$, which is the hypothesis, supplied twice.  Symmetry is the
exchange of the two conjuncts, the definition being symmetric under swapping
$a$ and $b$.
\end{proof}

\begin{designnote}[Why the tolerance exceeds one]\label{Order:note:delta}
The hypothesis $1 < \delta$ in $\name{within\_refl}$ is not a convenience; it
is forced.  Reflexivity requires $\name{ratio}(a,a) < \delta$, and
$\name{ratio}(a,a)$ is exactly $1$, so $\name{Within}(\delta,a,a)$ holds if and
only if $1 < \delta$.  More generally, $\name{ratio}(a,b)$ and
$\name{ratio}(b,a)$ are mutually inverse, so at least one is $\geq 1$ and a
tolerance $\delta \leq 1$ admits no pair at all.

So $\delta > 1$ is the multiplicative reading of $\varepsilon > 0$, and
$\delta = 1$ is the degenerate tolerance rather than the exact one --- the
opposite of what the additive analogy suggests.  The module header says as
much.  Only the sufficient direction is a declaration; that $1 < \delta$ is
also necessary is not stated in the file.
\end{designnote}

\begin{remark}[No triangle inequality]\label{Order:rem:no-triangle}
$\name{Within}$ is reflexive for admissible tolerances and symmetric, and
that is all.  The remaining ingredient of a metric-like notion --- the
multiplicative triangle inequality, that $\name{Within}(\delta,a,b)$ and
$\name{Within}(\varepsilon,b,c)$ give
$\name{Within}(\delta\varepsilon, a, c)$ --- is not proved here, and neither
is any statement combining two annuli.  Since Cauchy convergence is what the
header advertises this notion is for, the composition rule is the thing to
look for in the modules that follow.
\end{remark}

\subsection*{Summary of the correspondence}

\begin{center}
\small
\begin{tabular}{@{}lll@{}}
\hline
Lean declaration & Kind & Here \\
\hline
\lean{PrimeRatio.multiset\_eval\_pos\_order} & private thm & Lem.~\ref{Order:lem:pos} \\
\lean{PrimeRatio.lt\_of\_same\_left}  & private thm & Lem.~\ref{Order:lem:left} \\
\lean{PrimeRatio.lt\_same\_left}      & theorem     & Lem.~\ref{Order:lem:left} \\
\lean{PrimeRatio.lt\_of\_same\_right} & private thm & Lem.~\ref{Order:lem:right} \\
\lean{PrimeRatio.lt\_same\_right}     & theorem     & Lem.~\ref{Order:lem:right} \\
\lean{PrimeRatio.lt\_iff\_of\_same}   & theorem     & Lem.~\ref{Order:lem:both} \\
\lean{PrimeRatio.lt\_irrefl\_ratio}   & theorem     & Lem.~\ref{Order:lem:irrefl-rep} \\
\lean{PrimeRatio.lt\_trans\_ratio}    & theorem     & Lem.~\ref{Order:lem:trans-rep} \\
\lean{MulRat.lt}                      & definition  & Def.~\ref{Order:def:lt} \\
\lean{LT MulRat}                      & instance    & Def.~\ref{Order:def:lt} \\
\lean{MulRat.lt\_irrefl}              & theorem     & Thm.~\ref{Order:thm:strict} \\
\lean{MulRat.lt\_trans}               & theorem     & Thm.~\ref{Order:thm:strict} \\
\lean{MulRat.ratio}                   & definition  & Def.~\ref{Order:def:ratio} \\
\lean{MulRat.ratio\_self}             & simp thm    & Lem.~\ref{Order:lem:ratio-self} \\
\lean{MulRat.Within}                  & definition  & Def.~\ref{Order:def:within} \\
\lean{MulRat.within\_refl}            & theorem     & Lem.~\ref{Order:lem:within} \\
\lean{MulRat.within\_symm}            & theorem     & Lem.~\ref{Order:lem:within} \\
\hline
\end{tabular}
\end{center}

\noindent
Every declaration of \texttt{PrimeTensor/Order.lean} appears above.  The file
contains no \lean{sorry}, no axiom of its own, and no unproved named
proposition; it does use $\name{propext}$, as
Remark~\ref{Order:rem:propext} explains.  The necessity half of
Design note~\ref{Order:note:delta} is stated for the reader and is not a
declaration of the file.
